\par\noindent\textbf{(1)} 傅里叶定律 \(q=-kT'(x)\)；稳态无热源给 \(T''(x)=0\)，故 \(T(x)=T_h-(q/k)x\)。右面温度 \(T_s=T(L)=T_h-qL/k\)，温度沿板线性，但由于辐射边界 \(q\) 未知。
\par\noindent\textbf{(2)} 右面以面积热流形式满足 \(q=\epsilon\sigma(T_s^4-T_a^4)\)。代入 \(T_s=T_h-qL/k\) 可改写 \(k(T_h-T_s)/L=\epsilon\sigma(T_s^4-T_a^4)\)。对 \(T_h>T_a\)，方程在 \(T_a<T_s<T_h\) 有唯一解。
\par\noindent\textbf{(3)} 设 \(F(T_s)=k(T_h-T_s)/L-\epsilon\sigma(T_s^4-T_a^4)\)，在正温度区间 \(F'(T_s)=-k/L-4\epsilon\sigma T_s^3<0\)，且 \(F(T_a)>0,F(T_h)<0\)。故存在唯一交点，不会出现多个稳定稳态。
\par\noindent\textbf{(4)} 若 \(|T_h-T_a|\ll T_a\)，\(T_s^4-T_a^4\simeq4T_a^3(T_s-T_a)\)。设 \(h_r=4\epsilon\sigma T_a^3\)，则 \(q=(T_h-T_a)/[L/k+1/h_r]\)。这是导热热阻与线性化辐射热阻串联，不能在大温差时无条件使用。
